\chapter{Общий спектральный профиль отношения singlet--virtual}
\label{ch:v7-common-spectral-profile-singlet-virtual-ratio-gate}

\section{Какой коэффициент должен сократиться}

Совместный гессиан предыдущей главы зависит от двух безразмерных параметров
$a$ и $\gamma$. Проверим, превращает ли требование одного спектрального
профиля эту пару в вычисленное число. Для четырёхмерного действия
\begin{equation}
 \Tr\chi(D^2/\Lambda^2)
 \sim f_4\Lambda^4a_0+f_2\Lambda^2a_2+f_0a_4
 +\frac{f_{-2}}{\Lambda^2}a_6+\cdots
 \label{eq:v7-common-profile-heat-expansion}
\end{equation}
используем соглашение
\begin{equation}
 f_0=\chi(0),\qquad
 f_2=\int_0^\infty\chi(x)\,dx,\qquad
 f_{-2}=-\chi'(0).
 \label{eq:v7-common-profile-moments}
\end{equation}
Точные численные множители зависят от конечных следов, но распределение
уровней неизбежно: квадратичный Hodge-уровень приходит с $f_2\Lambda^2$,
квартика --- с $f_0$, а первый шестирёберный цикл --- с
$f_{-2}/\Lambda^2$ \cite{v7-common-profile-spectral-action,v7-common-profile-book}.

После измерения всех масс в Hodge-единицах остаётся безразмерный инвариант
профиля
\begin{equation}
 \boxed{R_\chi=\frac{f_2f_{-2}}{f_0^2}.}
 \label{eq:v7-common-profile-invariant}
\end{equation}
Конкретное $a$ получается умножением степени $R_\chi$ на вычислимые конечные
следовые коэффициенты полного product-оператора. Поэтому общий профиль
закрывает отношение только в том случае, если сам $R_\chi$ каноничен.

\section{Gaussian-зацепка проекта}

Для чистого теплового профиля
\begin{equation}
 \chi_t(x)=e^{-tx},\qquad t>0,
 \label{eq:v7-common-profile-pure-heat}
\end{equation}
имеем
\begin{equation}
 f_0=1,\qquad f_2=\frac1t,\qquad f_{-2}=t,
 \qquad R_{\chi_t}=1.
 \label{eq:v7-common-profile-gaussian-collapse}
\end{equation}
Тем самым Gaussian-линия ранних томов действительно даёт нетривиальную
подсказку: неизвестный тепловой масштаб сокращается из классического
отношения моментов. Это положительный частичный результат, а не право
назначить все конечные следовые множители единицами.

\section{Почему один профиль не равен одному параметру}

Рассмотрим столь же положительный профиль, являющийся смесью двух тепловых
масштабов:
\begin{equation}
 \chi(x)=w e^{-t_1x}+(1-w)e^{-t_2x},
 \qquad 0<w<1,\quad t_1,t_2>0.
 \label{eq:v7-common-profile-two-scale-mixture}
\end{equation}
Прямой расчёт даёт
\begin{equation}
 R_\chi
 =1+\frac{w(1-w)(t_1-t_2)^2}{t_1t_2}\ge1.
 \label{eq:v7-common-profile-mixture-ratio}
\end{equation}
Равенство достигается только при одном масштабе или на концах смеси. Уже
$w=1/2$, $t_1=1$, $t_2=4$ даёт
\begin{equation}
 R_\chi=\frac{25}{16}\ne1.
 \label{eq:v7-common-profile-mixture-benchmark}
\end{equation}
Следовательно, положительность, гладкость и даже представимость смесью
тепловых ядер не фиксируют нужное отношение. Чистый Gaussian является
дополнительным принципом экстремальности профиля.

С другой стороны, часто используемое усечение с плоским профилем в нуле
имеет
\begin{equation}
 \chi'(0)=0\quad\Longrightarrow\quad f_{-2}=0,
 \label{eq:v7-common-profile-flat-branch}
\end{equation}
и тогда классический коэффициент $a_6$ вместе с шестым циклическим словом
исчезает. Поэтому стандартное трёхмоментное усечение и ненулевой виртуальный
цикл нельзя использовать одновременно без явного неплоского профиля.

\section{Квантовая граница determinant-веса}

Даже выбор чистого Gaussian решает только классическое отношение
$a_2:a_4:a_6$. В четырёх измерениях интегрирование двух цветных мостов даёт
\begin{equation}
 -3\kappa^2|pq|^2
 \int\frac{d^4k}{(2\pi)^4}
 \frac1{(k^2+M_a^2)(k^2+M_b^2)},
 \label{eq:v7-common-profile-loop-integral}
\end{equation}
а этот интеграл логарифмически расходится. Перенормированный коэффициент
имеет структуру
\begin{equation}
 \lambda_{pq}^{\rm ren}(\nu)
 =\lambda_{pq}^{\rm ct}(\nu)
 +\lambda_{pq}^{\rm loop}(\nu;M_a,M_b,\kappa).
 \label{eq:v7-common-profile-renormalized-vertex}
\end{equation}
Масштаб согласования и конечная часть контрчлена не определяются голым
профилем $\chi$. Это стандартная граница однопетлевого исключения тяжёлых
полей \cite{v7-common-profile-hlm,v7-common-profile-dittmaier}. Поэтому
физический $\gamma$ совместного потенциала не выводится только из
$R_\chi=1$.

\begin{theorem}[Gaussian-сокращение и общий профильный no-go]
Чистый тепловой профиль $e^{-tx}$ точно даёт $R_\chi=1$ независимо от
$t$ и тем самым сокращает одну классическую свободу отношения Hodge- и
циклического уровней. Однако общий положительный профиль оставляет
$R_\chi$ переменным, а плоский в нуле профиль уничтожает шестой цикл.
Кроме того, четырёхмерный determinant-вес $\gamma$ требует независимого
перенормировочного условия. Следовательно, требование ``один спектральный
профиль'' не фиксирует физическое отношение $(a,\gamma)$.
\end{theorem}

\begin{nogocriterion}[Запрет считать функцию одним числом]
Нельзя выводить $a$ и $\gamma$ из факта существования общей функции
отсечения без фиксации её формы, коэффициента полного $a_6$, квантовой меры
и перенормировочного условия. Нельзя также одновременно пользоваться
плоским трёхмоментным усечением и ненулевым шестым циклическим членом.
Gaussian-профиль допускается только как явно названная дополнительная
гипотеза, а не как следствие спектрального принципа.
\end{nogocriterion}

\begin{protocol}
\begin{enumerate}
 \item Hodge-квартика: \textbf{уровень $f_0a_4$};
 \item Hodge-масса: \textbf{уровень $f_2\Lambda^2a_2$};
 \item шестой цикл: \textbf{уровень $f_{-2}a_6/\Lambda^2$};
 \item профильный инвариант: \textbf{$R_\chi=f_2f_{-2}/f_0^2$};
 \item чистый Gaussian: \textbf{$R_\chi=1$};
 \item положительная двухмасштабная смесь: \textbf{$R_\chi>1$ в общем случае};
 \item плоский профиль в нуле: \textbf{$f_{-2}=0$, цикл исчезает};
 \item конечный коэффициент полного $a_6$: \textbf{ещё не вычислен};
 \item четырёхмерный determinant-контрчлен: \textbf{необходим};
 \item $\gamma$ фиксирован профилем: \textbf{нет};
 \item коэффициент-свободное отношение $(a,\gamma)$: \textbf{не получено};
 \item итог: \textbf{Gaussian-классический частичный проход и общий no-go};
 \item следующий гейт: \textbf{полный product-$a_6$ коэффициент цикла}.
\end{enumerate}
\end{protocol}

Машинный сертификат сохранён в
\path{s2t_v7_common_spectral_profile_singlet_virtual_ratio_gate_results.json}.

\begin{thebibliography}{9}
\bibitem{v7-common-profile-spectral-action}
A.~H.~Chamseddine, A.~Connes,
\emph{The Spectral Action Principle}, Communications in Mathematical
Physics 186 (1997), 731--750; arXiv:hep-th/9606001.
\bibitem{v7-common-profile-book}
M.~Eckstein, B.~Iochum,
\emph{Spectral Action in Noncommutative Geometry}, SpringerBriefs in
Mathematical Physics 27 (2018); arXiv:1902.05306.
\bibitem{v7-common-profile-hlm}
B.~Henning, X.~Lu, H.~Murayama,
\emph{One-loop Matching and Running with Covariant Derivative Expansion},
Journal of High Energy Physics 01 (2018), 123; arXiv:1604.01019.
\bibitem{v7-common-profile-dittmaier}
S.~Dittmaier, S.~Schuhmacher, M.~Stahlhofen,
\emph{Integrating out heavy fields in the path integral using the
background-field method}, European Physical Journal C 81 (2021), 826;
arXiv:2102.12020.
\end{thebibliography}